% !TEX root = main.tex
% ACRONIMI — usare \gls{SIGLA} nel testo per la prima occorrenza espansa,
% le successive mostreranno solo la sigla. \glspl{SIGLA} per il plurale.
%
% Sintassi base:
%   \newacronym{SIGLA}{SIGLA}{forma estesa}
%
% Con plurali personalizzati:
%   \newacronym[longplural={forme estese}, shortplural={SIGLEs}]{SIGLA}{SIGLA}{forma estesa}
%
% NOTA: il glossary.tex originale della triennale conteneva ~26 acronimi attivi
% e ~140 commentati pronti all'uso, coprendo: 5G, GDPR, IoT, networking, cloud,
% sicurezza, machine learning, virtualizzazione. Recuperarli dalla cronologia git
% se si lavora in un dominio simile.

% Esempi:
% \newacronym{API}{API}{Application Programming Interface}
% \newacronym[longplural={Non-Functional Properties}, shortplural={NFPs}]{PNF}{PNF}{Proprietà Non Funzionale}
